\documentclass[10pt]{beamer}
\usetheme[
%%% options passed to the outer theme
%    progressstyle=fixedCircCnt,   %either fixedCircCnt, movCircCnt, or corner
%    rotationcw,          % change the rotation direction from counter-clockwise to clockwise
%    shownavsym          % show the navigation symbols
  ]{AAUsimple}
  
% If you want to change the colors of the various elements in the theme, edit and uncomment the following lines
% Change the bar and sidebar colors:
\setbeamercolor{AAUsimple}{fg=orange!20,bg=orange}
\setbeamercolor{sidebar}{bg=orange!20}
% Change the color of the structural elements:
\setbeamercolor{structure}{fg=black}
% Change the frame title text color:
\setbeamercolor{frametitle}{fg=white}
% Change the normal text color background:
%\setbeamercolor{normal text}{fg=black,bg=gray!10}
% ... and you can of course change a lot more - see the beamer user manual.

\usepackage[utf8]{inputenc}
\usepackage[spanish]{babel}
\usepackage[T1]{fontenc}
% Or whatever. Note that the encoding and the font should match. If T1
% does not look nice, try deleting the line with the fontenc.
\usepackage{helvet}
\usepackage[labelformat=empty]{caption}

% colored hyperlinks
\newcommand{\chref}[2]{%
  \href{#1}{{\usebeamercolor[bg]{AAUsimple}#2}}%
}

\title{Proyecto Grupal: Librería Para Generación de Horarios}
\subtitle{}  % could also be a conference name
\date{\today}

\author{
Emilio Rojas - B15680\\
  Juan Sánchez - B16068\\
  Marco Montero - A94000\\
  \href{mailto:}{{\tt}}
}

% - Give the names in the same order as they appear in the paper.
% - Use the \inst{?} command only if the authors have different
%   affiliation. See the beamer manual for an example

\institute[
%  {\includegraphics[scale=0.2]{aau_segl}}\\ %insert a company, department or university logo
  IE0217 - Estructuras de datos abstractas y algoritmos\\
  Escuela de Ingenieria Eléctrica\\
  Universidad de Costa Rica
] % optional - is placed in the bottom of the sidebar on every slide
{% is placed on the bottom of the title page
  IE0217 - Estructuras de datos abstractas y algoritmos\\\
  Escuela de Ingenieria Eléctrica\\
  Universidad de Costa Rica
  
  %there must be an empty line above this line - otherwise some unwanted space is added between the university and the country (I do not know why;( )
}

% specify a logo on the titlepage (you can specify additional logos an include them in 
% institute command below
\pgfdeclareimage[height=1.5cm]{titlepagelogo}{AAUgraphics/aau_logo_new} % placed on the title page
%\pgfdeclareimage[height=1.5cm]{titlepagelogo2}{AAUgraphics/aau_logo_new} % placed on the title page
\titlegraphic{% is placed on the bottom of the title page
  \pgfuseimage{titlepagelogo}
%  \hspace{1cm}\pgfuseimage{titlepagelogo2}
}

\begin{document}
% the titlepage
\aauwavesbg

\begin{frame}[plain,noframenumbering] % the plain option removes the header from the title page
\titlepage
\end{frame}
%%%%%%%%%%%%%%%%


%\begin{frame}{Índice}{}
%\tableofcontents
%
%\begin{figure}[!h]
%\begin{flushright}
%\includegraphics[width=.65\textwidth]{./AAUgraphics/data.jpg}
%%\caption{ \scriptsize } 
%\end{flushright}
%\end{figure}
%\end{frame}
%%%%%%%%%%%%%%%%







%%%%%%%%%%%%%%


%%%%%%%
\section{Introducción}

\begin{frame}{Texto ...}{}
\begin{block}{}
  \begin{itemize}
  \item bullets... Porque es importante, vision general, objetivos y alcvance tentativo (realista ojalá xD)
  \end{itemize}

%para colocar una figura...descomente esto!
%\begin{figure}[!h]
%\centering
%\includegraphics[width=0.9\textwidth]{./AAUgraphics/arbol.png}
%\caption{ \scriptsize } 
%\end{figure}

\end{block}
\end{frame}




\section{Implementación}
% motivation for creating this theme
\begin{frame}{Como lo vamos a hacer}{}
\begin{block}{}

\begin{itemize}
\item Más bullets... Planteamiento de la solución. Descrpción del diagrama de flujo del algoritmo, clases y métodos. explicar el encapsulamiento de clases que planeamos utilizar.
\end{itemize}
  
\end{block}
\end{frame}





\section{Características}
% motivation for creating this theme
\begin{frame}{Características}{}
\begin{block}{}

\begin{itemize}

\item Hablar de mySQL y porque lo utilizamos. También se puede mencionar la librería de c++ para accesar los datos y demás cosas relacionadfas con las bases de datos (estructura y keys importantes)

\end{itemize}

%Imagen
%\begin{figure}[!h]
%\centering
%\includegraphics[width=0.80\textwidth]{./AAUgraphics/insertion.png}
%\caption{ \scriptsize } 
%\end{figure}
 
\end{block}
\end{frame}




\section{En conclusión y logros hasta el momento }
\begin{frame}{Texto}{}
\begin{block}{}

\begin{itemize}
\item Problemas fundamentales (¿explicar cómo vamos a asignar a los profes a los cursos? (y por qué elegimos ese método)), ¿Qué tenemos hasta ahora?
\end{itemize}

%Imagen
%\begin{figure}[!h]
%\begin{flushright}
%\includegraphics[width=0.60\textwidth]{./AAUgraphics/node.png}
%\end{flushright}
%\end{figure}
  
\end{block}
\end{frame}






\begin{frame}{Bibliografía}{}
\begin{block}{}
\begin{itemize}

\item Robert Sedgewick. \textbf{Left-Leaning Red-Black Trees}. Princeton University. 2008

\item CS 660: Combinational Algorithms. \textbf{(a,b)- Trees}. San Diego University
\end{itemize}
\end{block}
\end{frame}


\begin{frame}{}{}
\begin{block}{}

\begin{figure}[!h]
\begin{center}
\includegraphics[width=.8\textwidth]{./AAUgraphics/gracias.jpg}
%\caption{ \scriptsize } 
\end{center}
\end{figure}

\end{block}
\end{frame}




%\begin{frame}{}{}
%\begin{block}{}
%
%\begin{figure}[!h]
%
%\begin{center}
%\includegraphics[width=.8\textwidth]{./AAUgraphics/gracias.jpg}
%\caption{ \scriptsize } 
%\end{center}
%\end{figure}
%
%\end{block}
%\end{frame}

\begin{frame}{Preguntas}{}
\begin{block}{}

 \begin{figure}[!h]
 \begin{flushleft}
\includegraphics[width=.25\textwidth]{./AAUgraphics/preguntas1.jpg}
%\caption{ \scriptsize }
\end{flushleft} 
\end{figure}

\begin{figure}[!h]
\begin{flushright}
\includegraphics[width=.25\textwidth]{./AAUgraphics/preguntas2.jpg}
%\caption{ \scriptsize } 
\end{flushright}
\end{figure}

\end{block}
\end{frame}




\end{document}
%%%%% HASTA AQUí!!!! ------FIN------- %%%%%%%%%%%%%%%%%%%%%%
